% Source: https://artofproblemsolving.com/community/c3h2655593_a_cute_problem_
% Stable identity: aops-topic-2655593-q01
The first factor is $1+2+3=6$. In the double sum, for each $n\ge2$ there are exactly $n-1$ pairs of positive integers $(j,k)$ with $j+k=n$. All terms are positive, so we may regroup:
\[\sum_{j=1}^{\infty}\sum_{k=1}^{\infty}\dfrac{1}{(j+k)^3}=\sum_{n=2}^{\infty}\dfrac{n-1}{n^3}=\sum_{n=2}^{\infty}\left(\dfrac1{n^2}-\dfrac1{n^3}\right).\]
Adding $\displaystyle\sum_{i=1}^{\infty}\dfrac1{i^3}=1+\sum_{n=2}^{\infty}\dfrac1{n^3}$, the cubes cancel:
\[\sum_{i=1}^{\infty}\dfrac1{i^3}+\sum_{j,k\ge1}\dfrac{1}{(j+k)^3}=1+\sum_{n=2}^{\infty}\dfrac1{n^2}=\dfrac{\pi^2}6.\]
So the value is $\sqrt{6\cdot\dfrac{\pi^2}{6}}=\boxed{\pi}$.
